\chapter{Anfangsabschnitte, Schnitte und Sprünge}
\label{b15guide:ch:sections}

\section{Anfangs- und Endabschnitte}

\begin{definition}[Anfangsabschnitt und Endabschnitt]
\label{b15guide:def:sections}
$L\subseteq A$ heißt \emph{Anfangsabschnitt}, wenn aus
$x\in L$ und $y\leq x$ stets $y\in L$ folgt.
$U\subseteq A$ heißt \emph{Endabschnitt}, wenn aus
$x\in U$ und $x\leq y$ stets $y\in U$ folgt.
Die Mengen $\varnothing$ und $A$ sind zugelassen.
\end{definition}

\begin{theorem}[Trennung eines Anfangsabschnitts vom Rest]
\label{b15guide:thm:initial-separation}
Für $L\subseteq A$ sind äquivalent:
\[
\begin{aligned}
 &L\text{ ist ein Anfangsabschnitt},\\
 &\forall l\in L\,\forall u\in A\setminus L\;(l<u).
\end{aligned}
\]
Dual ist $U$ genau dann ein Endabschnitt, wenn
jedes Element von $A\setminus U$ kleiner als jedes Element von $U$ ist.
\end{theorem}

\begin{theorem}[Anfangsabschnitte sind durch Inklusion total geordnet]
\label{b15guide:thm:nested}
Für zwei Anfangsabschnitte $L,K$ gilt stets
\[
 L\subseteq K\quad\text{oder}\quad K\subseteq L.
\]
Dasselbe gilt für zwei Endabschnitte. Insbesondere ist
der Schnitt zweier Anfangsabschnitte einer der beiden Abschnitte,
ebenso ihre Vereinigung; die entsprechende Aussage gilt für Endabschnitte.
\end{theorem}

\begin{theorem}[Abschnittskriterium für Totalität]
\label{b15guide:thm:initial-criterion}
Für eine halbgeordnete Menge $(P,\preceq)$ sind äquivalent:
\begin{enumerate}
 \item $\preceq$ ist total.
 \item Je zwei Anfangsabschnitte von $P$ sind durch Inklusion vergleichbar.
 \item Je zwei Mengen $\{z\in P\mid z\preceq x\}$ und
       $\{z\in P\mid z\preceq y\}$ sind durch Inklusion vergleichbar.
\end{enumerate}
Dual gelten dieselben Äquivalenzen mit Endabschnitten
und den Mengen $\{z\in P\mid x\preceq z\}$.
\end{theorem}

\begin{theorem}[Strikte Strahlen bestimmen den Vergleich]
\label{b15guide:thm:rays}
Für $a,b\in A$ gilt
\[
\begin{aligned}
 a\leq b&\Longleftrightarrow A_{<a}\subseteq A_{<b},&
 a<b&\Longleftrightarrow A_{<a}\subsetneq A_{<b},\\
 a\leq b&\Longleftrightarrow A_{>b}\subseteq A_{>a},&
 a<b&\Longleftrightarrow A_{>b}\subsetneq A_{>a}.
\end{aligned}
\]
Insbesondere bestimmt $A_{<a}$ das Element $a$ eindeutig;
dasselbe gilt für $A_{>a}$.
\end{theorem}

\section{Schnitte ohne Vollständigkeitsannahme}

\begin{definition}[Schnitt und seine Endpunkte]
\label{b15guide:def:cut}
Ein \emph{Schnitt} von $A$ ist hier ein Paar $(L,U)$ mit
\[
 A=L\mathbin{\dot\cup}U,\qquad
 \forall l\in L\,\forall u\in U\;(l<u).
\]
Leere Seiten sind erlaubt. Der Schnitt heißt \emph{echt},
wenn beide Seiten nichtleer sind.
Ein größtes Element von $L$ heißt \emph{linker Endpunkt},
ein kleinstes Element von $U$ \emph{rechter Endpunkt} des Schnitts.
Diese Konvention verlangt weder einen fehlenden linken Endpunkt
noch die Existenz eines Elements an der Schnittstelle.
\end{definition}

\begin{theorem}[Schnitte und Anfangsabschnitte]
\label{b15guide:thm:cut-initial}
Die Schnitte von $A$ sind genau die Paare
\[
 (L,A\setminus L)
 \quad\text{mit }L\text{ einem Anfangsabschnitt}.
\]
Ihre linken Seiten sind unter Inklusion total geordnet;
aus $L\subseteq L'$ folgt für die zugehörigen rechten Seiten
$U'\subseteq U$.
\end{theorem}

\begin{theorem}[Endpunktdarstellung eines Schnitts]
\label{b15guide:thm:cut-endpoint}
Für einen Schnitt $(L,U)$ gilt:
\[
\begin{aligned}
 a=\max L&\Longleftrightarrow
          (L=A_{\leq a}\ \land\ U=A_{>a}),\\
 b=\min U&\Longleftrightarrow
          (L=A_{<b}\ \land\ U=A_{\geq b}).
\end{aligned}
\]
Jede Gleichheit mit $\min$ oder $\max$ schließt die Existenz
des entsprechenden Endpunkts ein.
\end{theorem}

\begin{theorem}[Die vier Endpunktfälle eines echten Schnitts]
\label{b15guide:thm:cut-types}
Jeder echte Schnitt $(L,U)$ gehört zu genau einem der Fälle:
\begin{enumerate}
 \item $L$ hat ein Maximum $a$, und $U$ hat ein Minimum $b$.
       Dann ist $a<b$ und $(a,b)=\varnothing$.
 \item $L$ hat ein Maximum $a$, und $U$ hat kein Minimum.
       Dann ist $(L,U)=(A_{\leq a},A_{>a})$.
 \item $L$ hat kein Maximum, und $U$ hat ein Minimum $b$.
       Dann ist $(L,U)=(A_{<b},A_{\geq b})$.
 \item $L$ hat kein Maximum, und $U$ hat kein Minimum.
\end{enumerate}
Im vierten Fall sprechen wir hier von einer \emph{Lücke}.
Totalität allein schließt weder Sprünge noch Lücken aus.
\end{theorem}

\section{Unmittelbare Nachbarn}

\begin{definition}[Unmittelbarer Nachfolger und Vorgänger]
\label{b15guide:def:neighbor}
Für $a,b\in A$ heißt $b$ \emph{unmittelbarer Nachfolger} von $a$,
und $a$ \emph{unmittelbarer Vorgänger} von $b$, wenn
\[
 a<b\quad\text{und}\quad
 \neg\exists x\in A\;(a<x<b).
\]
Das Paar $(a,b)$ heißt dann ein \emph{Sprung}.
\end{definition}

\begin{theorem}[Charakterisierung und Eindeutigkeit der Nachbarn]
\label{b15guide:thm:neighbor}
Für $a,b\in A$ sind äquivalent:
\begin{enumerate}
 \item $b$ ist unmittelbarer Nachfolger von $a$.
 \item $b=\min A_{>a}$.
 \item $a=\max A_{<b}$.
 \item $a<b$ und $A_{\leq a}=A_{<b}$.
\end{enumerate}
Die Minimums- und Maximumsaussagen schließen deren Existenz ein.
Jedes Element hat höchstens einen unmittelbaren Nachfolger
und höchstens einen unmittelbaren Vorgänger.
\end{theorem}

\begin{theorem}[Nachbarn in endlichen Teilmengen]
\label{b15guide:thm:finite-neighbors}
Sei $S=\{s_1,\ldots,s_n\}\subseteq A$ nichtleer und endlich
mit $s_1<\cdots<s_n$. In der auf $S$ eingeschränkten Ordnung
sind genau die Paare $(s_i,s_{i+1})$, $1\leq i<n$, Sprünge.
Insbesondere hat jedes Element von $S$ außer $\max S$ einen
unmittelbaren Nachfolger in $S$, und jedes außer $\min S$
einen unmittelbaren Vorgänger in $S$.
Diese Nachbarschaft wird relativ zu $S$ behauptet.
\end{theorem}
